\documentclass[../main.tex]{subfiles}
\graphicspath{{\subfix{../images/}}}

\begin{document}
\section{Discussion}
% NOTE: Each para should start with what you found and then go into what other people found and then how it is the same/different!! 
% para 1: summarize findings
% para 7: multiplicative effects 
% para 5: summary of previous PO studies + health and contextualize our study 
% para 8: vulnerable pops + copd
% para 6: wfs main effects and results and contextualize study within the context of this lit
% para 3: describe PO exposure measurement
% para 4: transition from exp measurement is diff to also our exp is different... summary of previous PSPS studies + health and again contextualize our study 
% para 9: limitations 
% para 10: conc

% para 1: summarize findings
California's wildfire frequency and severity have increased rapidly.\citep{dennison2014large, li2021spatial, cal_fire_2022} In tandem, utilities have increasingly initiated preemptive power shutoffs to reduce fire risk and spread.\citep{huang2023review} Using California data from 2013--2019, we found a positive association between severe magnitude PSPS events and same-day all-cause respiratory and COPD-specific, but not cardiovascular or psychiatric, emergency department visits and hospital admissions. Wildfire \PM exposure alone was also only associated with same-day all-cause respiratory and COPD-specific emergency department visits and hospital admissions. The association between PSPS events and respiratory endpoints was stronger in the presence of wildfire \PM exposure, especially for COPD. 

% para 7: multiplicative effects 
While both power outages and wildfire smoke are associated with adverse health outcomes, the effects of these two exposures together are understudied, despite a rise in their co-occurrence.\citep{do2025spatiotemporal} This co-occurring exposure may have more than additive health effects. Several studies found that power outage co-occurring with extreme heat, storms, wind events, and floods have been shown to exacerbate health effects.\citep{xu2025power, yamasaki2024heat, lin2011health, lin2021immediate, lin2024joint} We found evidence of both additive and multiplicative interaction between PSPS events and wildfire \PM on respiratory healthcare utilization, particularly for COPD. The higher additive interaction for COPD was driven by a stronger relationship between PSPS events and COPD encounters than overall respiratory encounters. The multiplicative interaction between PSPS events and wildfire \PM exposure was of similar magnitude for COPD and overall respiratory encounters. These findings suggest that climate-driven hazards may be risk multipliers for power outage events.

% para 5: summary of previous PO studies + health and contextualize our study 
Limited large-scale epidemiological studies have evaluated the relationship between any type of power outage and health.\citep{casey2020power, northrop2024power, lin2021immediate, mcbrien2025, molinari2017s, zhang2020power}  Generally, the existing literature finds moderately elevated health risks following power outages, with some evidence that the highest risk is concentrated among those dependent on DME.\citep{casey2020power, northrop2024power, lin2021immediate, mcbrien2025, molinari2017s, andresen2023understanding} In the days following a power outage, most, but not all,\citep{lin2021immediate, do2025impact} studies in New York State and Florida have found increased risk of emergency department visits or hospitalizations for respiratory illness,\citep{dominianni2018health, lin2024joint, deng2022independent} COPD,\citep{zhang2020power, lin2024joint} cardiovascular disease,\citep{deng2022independent} food and water-borne disease,\citep{deng2022independent} injury,\citep{northrop2024power} renal disease,\citep{dominianni2018health} and carbon monoxide poisoning.\citep{northrop2024power} Our results for severe magnitude PSPS events and respiratory encounters align with the majority of the existing literature, as do our null results for cardiovascular disease,\citep{lin2021immediate, do2025impact} suggesting respiratory disease may be more sensitive to power outages than cardiovascular disease. We found no association with psychiatric encounters, despite participants in two qualitative studies reporting increases in worry, anxiety, and stress following a power outage.\citep{ibrahim2016erratic, moreno2019community} Although people report worry around PSPS events,\citep{WONGPARODI2022102495} their symptoms may not rise to a level warranting emergency psychiatric medical care. 

The population health literature around PSPS events remains sparse.\citep{WONGPARODI2022102495, harding2025a} Motor vehicle crashes appear largely unchanged during these events.\citep{harding2025a} However, because numerous states now initiate PSPS events, future work must evaluate their impacts on health to inform cost-benefit decision making. 

% para 8: vulnerable pops + copd
Prior studies have found individuals reliant on electricity-dependent durable medical equipment, older adults, and those with vs. without pre-existing conditions are at higher risk following power outage exposure.\citep{molinari2017s, harding2024association, sheridan2021individual, deng2022independent, xu2025power, casey2020power} PSPS events likely pose particular danger to COPD patients, given their reliance on DME, sensitivity to temperature and humidity changes, and frequent comorbidities.\citep{casey2020power} Our findings are consistent, with the strongest associations between PSPS events and emergency department visits and hospitalizations being for COPD encounters. A large portion (\> 10\%) of those with COPD rely on long-term oxygen therapy.\citep{hess2023oxygen, nishi2015oxygen} In New York State, power outages were associated with increased risk of COPD-related hospitalization, and power outages appear to mediate the association between severe weather and COPD healthcare utilization.\citep{zhang2020power, qu2021power}

% para 6: wfs main effects and results and contextualize study within the context of this lit
Short-term wildfire \PM exposure has been associated with increased risk of respiratory hospitalizations and ED visits, but not consistently with cardiovascular encounters.\citep{gould2024health} Studies have also identified associations between short-term wildfire \PM exposure and COPD hospital encounters.\citep{sutherland2005wildfire, reid2019wildfire, wilgus2024clearing} Results are mixed for mental health-related healthcare utilization, with some evidence of same-day increases.\citep{cleland2022short, wettstein2024psychotropic, jung2025fine, chen2023emergency} Our study generally agrees with this existing literature. 

% para 3: describe PO exposure measurement
The majority of epidemiological studies on power outage have constructed binary power outage metrics. These metrics defined a power outage based on a certain percentage of individuals lacking power (e.g., 5\%) in a specific geographic unit.\citep{northrop2024power, do2025impact, lin2021immediate, lin2024joint, zhang2020power, deng2022independent} Others considered everyone in an affected area as exposed (e.g., New York during the 2003 blackout).\citep{anderson2011impacts, dominianni2018health} One study used outage data at specific skilled nursing facilities.\citep{skarha2021association} 

% para 4: transition from exp measurement is diff to also our exp is different... summary of previous PSPS studies + health and again contextualize our study 
Our power outage exposure assessment deviated in several key ways from previous studies. First, we only evaluated PSPS events, which are always planned, proactive cuts to power. The decision to initiate a PSPS event reflects a balancing of the societal risks stemming from wildfires and the costs of the PSPS event itself.\citep{huang2023review} Power utility companies are required to inform customers in advance and take precautions to mitigate the impacts of the outage, such as providing access to Community Resource Centers where individuals can charge devices and access WiFi.\citep{harding2025a} Historically, PSPS events last longer than other types of power outages. Importantly, unlike previous studies, our primary exposure metric was based on counts of customers without power, rather than duration. This is because nearly all PSPS events lasted longer than 8 hours, the cutoff generally used in previous health studies.\citep{do2025spatiotemporal, northrop2024power} 

% para 9: limitations 
\subsection{Limitations}
This study had several limitations. The exposure assessment leveraged power utility data at the circuit level, meaning we did not have access to individual-level exposure data. We aggregated circuit-level data to the zip code-level, which resulted in some exposure misclassification. Further, when estimating exposure data at the residential zip code-level, we could not ensure that an individual was at home when a PSPS event occurred. This limitation also applied to wildfire \PM exposure ascertainment. 

Second, the average duration of PSPS events spanned multiple days, and we examined only same-day exposure. Future studies should examine lagged effects of PSPS events and wildfire \PM exposure. 

Third, we only captured individuals 20 years and older who went to an emergency department or hospital. Individuals who seek healthcare may systematically differ from those who do not. Additionally, PSPS events may result in delayed care-seeking behavior, and while we only assessed same-day visits, future research should consider more extended periods. Because effects in children may differ, future studies should include their healthcare information. 

Fourth, we grouped many disease healthcare utilization together into four categories (respiratory, COPD, psychiatric, and cardiovascular). This may explain our null results for overall psychiatric endpoints. There is a wide discrepancy in care-seeking behavior for some psychiatric endpoints, such as depression compared to substance abuse. Further, we grouped emergency department visits with hospital admissions, however the decision to seek care in these settings differs. Patients or families decide to seek care in the case of emergency department encounters, whereas an admitting physician decides to admit a patient. Future studies may wish to disaggregate the types and causes of encounters.

Finally, in an observational study, unobserved confounding may remain. Our case-crossover design effectively handles time-invariant confounding, but short-term time-varying confounders may remain. We adjusted models for daily ambient temperature, but there may be residual confounding.

% para 9: conc
\subsection{Conclusion}
In the US, the prevalence of wildfires and co-occurring power outages is increasing\citep{iglesias2022us, climatecentral2024weather, do2025spatiotemporal}. Since 2022, states other than California have also implemented PSPS events, including Washington, Oregon, Idaho, and Hawaii\citep{PSEHealthyEnergy_PSPS} in response to growing wildfire risk. Our study provides the first analysis of these co-occurring hazards related to healthcare utilization. Given the increasing frequency of PSPS events, research is essential to understand their public health implications both independently and alongside wildfire \PM, ensuring these mitigation strategies reduce rather than elevate risks to human health.

\end{document}